\documentclass[11pt,a4paper]{article}

\usepackage[utf8]{inputenc}
\usepackage[T1]{fontenc}
\usepackage{amsmath,amssymb,amsfonts}
\usepackage{booktabs}
\usepackage{graphicx}
\usepackage{hyperref}
\usepackage{xcolor}
\usepackage{tikz}
\usetikzlibrary{positioning,arrows.meta,calc,fit,backgrounds}
\usepackage[margin=2.5cm]{geometry}
\usepackage{caption}
\usepackage{subcaption}
\usepackage{natbib}
\bibliographystyle{plainnat}

\newcommand{\R}{\mathbb{R}}
\newcommand{\E}{\mathbb{E}}
\newcommand{\bW}{\mathbf{W}}
\newcommand{\bb}{\mathbf{b}}
\newcommand{\bh}{\mathbf{h}}
\newcommand{\bx}{\mathbf{x}}
\newcommand{\ba}{\mathbf{a}}
\newcommand{\bv}{\mathbf{v}}
\newcommand{\bo}{\mathbf{o}}
\newcommand{\relu}{\mathrm{GELU}}
\newcommand{\lnorm}{\mathrm{LN}}
\newcommand{\fc}{\mathrm{fc}}
\newcommand{\attnout}{\mathbf{a}_{\mathrm{attn}}}
\newcommand{\cosim}{\mathrm{cos}}

\title{How a One-Layer Transformer Implements\\
  Relation-Conditioned Affine Maps}
\author{}
\date{}

\begin{document}
\maketitle

\begin{abstract}
The linear representation hypothesis proposes that language models
implement relational knowledge via affine transformations in their
representation spaces. We provide a complete mechanistic account of how
a one-layer transformer learns relation-conditioned affine maps in a
controlled synthetic setting. The model achieves $R^2 \approx 0.98$ for
affine approximations of its MLP computation, with 99.7\% accuracy in
mapping subject representations to correct type centroids. The mechanism
decomposes into three steps: (1)~a single attention head injects a
relation-identity direction into the residual stream, superimposed on
and \emph{99\% orthogonal} to the subject embedding; (2)~each MLP
neuron's input weight vector independently projects onto both signals
--- the relation component gates the neuron via GELU, while the subject
component modulates activation magnitude; (3)~the active neurons' output
vectors form a relation-specific dictionary mapping subject variation to
object centroids. The learned orthogonality between relation and subject
signals is what allows a single weight matrix to implement two
different affine transformations without interference.
\end{abstract}

%======================================================================
\section{Introduction}
\label{sec:intro}
%======================================================================

The linear representation hypothesis
\citep{hernandez2024linearity,park2024linear} proposes that neural
language models encode relational knowledge through affine
transformations: given a subject representation~$\mathbf{s}$, the model
computes the corresponding object representation as $\hat{\bo} =
\bW\mathbf{s} + \bb$. Empirical evidence shows that such affine probes
can recover relational structure from intermediate representations with
high accuracy. However, the question of \emph{how} the network
implements these affine maps --- which components contribute, and
through what mechanism --- has received less attention.

We study this question in a controlled synthetic setting where
ground-truth relations are known exactly, allowing complete mechanistic
analysis. We train a one-layer GPT-style transformer on a task with
2025 entities organized into 45 types and two relations (R0 and R1),
each defined by shared attribute types. The model achieves 100\%
training accuracy and 90.6\% generalization to held-out (subject,
relation) pairs through pure type-level generalization.

Our main finding is a three-step mechanism by which the model
implements relation-conditioned affine maps:

\begin{enumerate}
\item \textbf{Superposition.} A single attention head reads the
  relation token and adds a relation-identity direction to the residual
  stream. This direction is 99\% orthogonal to the subject-embedding
  subspace, so the two signals coexist without interference.

\item \textbf{GELU gating.} Each MLP neuron's input weight vector
  projects independently onto the relation direction (for gating) and
  the subject subspace (for type-specific computation). The relation
  projection determines whether the neuron is above or below GELU's
  zero crossing, partitioning the 256 neurons into two
  relation-specific sub-networks (85 for R0, 82 for R1).

\item \textbf{Dictionary readout.} The active neurons' output weight
  columns form a learned, approximately isotropic dictionary. The
  relation-specific activation pattern selects which dictionary elements
  contribute, producing a relation-conditioned affine output.
\end{enumerate}

The result is that the entire MLP, for a given relation~$k$, is well
approximated by a single affine map $\mathrm{MLP}(\bx) \approx
\bW_k \bx + \bb_k$, with $R^2 \approx 0.99$. GELU's primary role is
not nonlinear computation but \emph{gating}: selecting which relation's
linear map to apply.

%======================================================================
\section{Experimental Setup}
\label{sec:setup}
%======================================================================

\paragraph{Entities and types.}
We define 2025 entities, each characterized by two attribute groups:
an \emph{f-type} and a \emph{g-type}, where each type is a 2-element
subset of $\{0,\dots,9\}$, giving $\binom{10}{2} = 45$ possible types
per group. Two relations are defined over these entities:
\begin{itemize}
\item \textbf{R0}: subject and object share the same f-type.
\item \textbf{R1}: subject and object share the same g-type.
\end{itemize}
Each (subject, relation) pair maps to 45 valid objects (those sharing
the target type). The task is next-token prediction: given context
$[\texttt{sep},\; r_k,\; \text{subj}_i]$, predict a valid object.

\paragraph{Architecture.}
We use a one-layer GPT-style causal transformer with hidden
dimension~64, 4 attention heads (head dimension~16), and an MLP with
4$\times$ expansion (256 neurons) and GELU activation. The model has
approximately 580{,}000 parameters. The architecture follows the
standard pre-norm residual pattern:
\begin{align}
\bh_0 &= \mathrm{tok\_emb}(\text{input}) + \mathrm{pos\_emb}(\text{positions}), \label{eq:h0}\\
\bh_0' &= \bh_0 + \mathrm{Attn}(\lnorm_1(\bh_0)), \label{eq:attn}\\
\bh_1 &= \bh_0' + \mathrm{MLP}(\lnorm_2(\bh_0')), \label{eq:mlp}\\
\text{logits} &= \mathrm{lm\_head}(\lnorm_f(\bh_1)), \label{eq:logits}
\end{align}
where $\mathrm{MLP}(\bx) = \fc_2(\relu(\fc_1(\bx)))$, with
$\fc_1: \R^{64} \to \R^{256}$ and $\fc_2: \R^{256} \to \R^{64}$.

\paragraph{Training.}
Training data consists of 3-token attribute sentences
($[\texttt{sep},\;\text{entity},\;\text{attribute}]$) and 4-token
relation triples
($[\texttt{sep},\; r_k,\; \text{subj}_i,\; \text{obj}_j]$), packed
into sequences. We train for 10 epochs with cosine learning rate
schedule (peak $3 \times 10^{-4}$, 200-step warmup). The model
achieves 100\% accuracy on trained (subject, relation) pairs and
90.6\% generalization to held-out pairs --- pure type-level
generalization without ever observing those specific pairs during
training.

%======================================================================
\section{The Affine Map Result}
\label{sec:affine}
%======================================================================

We assess whether the model's internal computation can be described by
a relation-conditioned affine transformation. For each relation~$k$,
we fit a least-squares affine map $\hat{\bo} = \bW_k \mathbf{s} +
\bb_k$ from subject hidden states~$\mathbf{s}$ to target type
centroids (the mean \texttt{lm\_head} weight row across all 45 valid
objects). We evaluate using Voronoi accuracy (whether the affine
prediction lands nearest to the correct centroid among all 45 type
centroids) and cosine similarity.

\paragraph{Residual stream analysis.}
We fit the affine map at four points in the residual stream, using a
50/50 train/test split over subjects (Table~\ref{tab:residual}). The
affine map achieves near-perfect Voronoi accuracy at \emph{every}
point, including the raw embedding~$\bh_0$ before any attention has
been applied.

\begin{table}[t]
\centering
\caption{Affine map accuracy at each residual stream point (v7
  checkpoint, 1012 train / 1013 test subjects). Target: type centroid.}
\label{tab:residual}
\begin{tabular}{lcccc}
\toprule
Residual point & R0 train & R0 test & R1 train & R1 test \\
\midrule
\multicolumn{5}{l}{\emph{Voronoi accuracy}} \\
$\bh_0$ (embedding) & 0.997 & 0.997 & 0.998 & 0.998 \\
$\bh_0 + \text{attn}$ & 0.996 & 0.996 & 0.998 & 0.998 \\
$\bh_0 + \text{attn} + \text{mlp}$ & 0.997 & 0.994 & 0.999 & 1.000 \\
$\lnorm_f(\bh_1)$ & 0.997 & 0.995 & 0.999 & 0.999 \\
\midrule
\multicolumn{5}{l}{\emph{Cosine similarity}} \\
$\bh_0$ (embedding) & 0.968 & 0.965 & 0.973 & 0.966 \\
$\bh_0 + \text{attn}$ & 0.968 & 0.965 & 0.972 & 0.966 \\
$\bh_0 + \text{attn} + \text{mlp}$ & 0.969 & 0.967 & 0.974 & 0.969 \\
$\lnorm_f(\bh_1)$ & 0.970 & 0.968 & 0.975 & 0.970 \\
\bottomrule
\end{tabular}
\end{table}

That the raw embedding $\bh_0 = \mathrm{tok\_emb}(\text{subj}_i) +
\mathrm{pos\_emb}(2)$ already supports 99.7\% Voronoi accuracy with
\emph{separate} affine maps for R0 and R1 is a key observation: the
64-dimensional embedding encodes both f-type and g-type information in
a form that is linearly accessible. Different affine maps for R0
vs.\ R1 simply read out the appropriate type dimension. The slight but
consistent increase in cosine similarity through the layers
($0.965 \to 0.970$) shows the MLP actively refines the representation
toward the correct type centroid.

\paragraph{MLP-internal affine fit.}
We also fit affine maps to the MLP's internal computation. For each
relation, we approximate the mapping from the 64-dimensional MLP input
to the 256-dimensional post-GELU activations (Table~\ref{tab:mlp_affine}).

\begin{table}[t]
\centering
\caption{Affine approximation of MLP computation, per relation.}
\label{tab:mlp_affine}
\begin{tabular}{lcc}
\toprule
Metric & R0 & R1 \\
\midrule
\multicolumn{3}{l}{\emph{MLP input $\to$ post-GELU activations}} \\
Overall $R^2$ & 0.980 & 0.976 \\
Per-entity cosine (mean) & 0.998 & 0.997 \\
Per-entity relative error & 0.066 & 0.070 \\
\midrule
\multicolumn{3}{l}{\emph{MLP input $\to$ MLP output (end-to-end)}} \\
Overall $R^2$ & 0.990 & 0.989 \\
Per-entity cosine (mean) & 0.997 & 0.996 \\
\midrule
\multicolumn{3}{l}{\emph{Comparison with $\fc_1$ weights}} \\
$\cosim(\text{fitted col},\; \fc_1\text{ col})$ & 0.000 & $-0.004$ \\
\bottomrule
\end{tabular}
\end{table}

The affine approximation explains 97.6--98.0\% of the variance in
post-GELU activations. The end-to-end fit (MLP input to MLP output)
is even better ($R^2 \approx 0.99$), because the $\fc_2$ projection
removes residual noise in the 256-dimensional activation space.
Importantly, the fitted affine map is \emph{not} a scaled version of
$\fc_1$: the cosine similarity between fitted columns and $\fc_1$
columns is approximately zero. The affine map absorbs the
relation-dependent gating into its parameters, learning an effective
linear map that accounts for which neurons are active and at what
operating point on the GELU curve.

%======================================================================
\section{Superposition in the Residual Stream}
\label{sec:superposition}
%======================================================================

How does a single MLP with one set of weights implement two different
affine maps --- one for each relation? The answer lies in the
structure of its input. At the subject position (position~2 in the
context $[\texttt{sep},\; r_k,\; \text{subj}_i]$), the MLP input is
\begin{equation}
\bx = \lnorm_2(\bh_0 + \attnout),
\label{eq:mlp_input}
\end{equation}
where $\bh_0 = \mathrm{tok\_emb}(\text{subj}_i) +
\mathrm{pos\_emb}(2)$ encodes subject identity and $\attnout$ is the
attention output at position~2. These two signals carry fundamentally
different information and are superimposed additively in the residual
stream.

%----------------------------------------------------------------------
\subsection{Attention output is a near-constant relation signal}
\label{sec:attn_signal}
%----------------------------------------------------------------------

We decompose the attention output at position~2 across all 2025
subjects under each relation. Of the total variance in $\attnout$,
\textbf{95.6\%} is between relations and only 4.4\% within. The
within-relation pairwise cosine similarity is 0.96--0.99, meaning the
attention output is nearly identical for all subjects queried with the
same relation. The cross-relation cosine is only 0.10. Projecting
each $\attnout$ onto the relation direction $\bar{\ba}_{\text{R0}} -
\bar{\ba}_{\text{R1}}$ yields a Cohen's $d$ of~27 --- effectively
zero overlap between R0 and R1. In summary:
\begin{equation}
\attnout \;\approx\; \bv_{\text{rel}} + \boldsymbol{\varepsilon}_{\text{subj}},
\label{eq:attn_decomp}
\end{equation}
where $\bv_{\text{rel}}$ is a constant vector determined by the
relation token and $\boldsymbol{\varepsilon}_{\text{subj}}$ is
negligible.

Decomposing by head reveals that a \textbf{single head} (head~1) is
responsible for 95.2\% of the relation direction. Head~1 attends
73--78\% to the relation token at position~1, producing outputs with
within-relation cosine 0.999 and R0/R1 cosine $-0.48$ (opposite
directions). Its output norm (0.26--0.31) is 2--4$\times$ larger than
the other heads. The remaining heads play minor roles: head~0 provides
a constant offset (attending primarily to \texttt{sep}), head~2
contributes 4.8\% of the relation direction, and head~3 carries the
most subject-specific variation but has the smallest norm.  Full
per-head statistics are in Appendix~\ref{app:attention}.

%----------------------------------------------------------------------
\subsection{How $\fc_1$ reads the superposition}
\label{sec:fc1_reads}
%----------------------------------------------------------------------

Each row of $\fc_1$ (the input weight vector for neuron~$j$) projects
the superposed MLP input:
\begin{equation}
\mathrm{pre}_j = \fc_{1,j} \cdot \bx + b_j
  = \underbrace{\fc_{1,j} \cdot \bar{\bx}_k + b_j}_{\text{relation component}}
  + \underbrace{\fc_{1,j} \cdot (\bx - \bar{\bx}_k)}_{\text{subject component}},
\label{eq:pre_decomp}
\end{equation}
where $\bar{\bx}_k$ is the mean MLP input for relation~$k$. The
relation component is constant for all subjects under the same
relation and determines the neuron's operating regime. The subject
component varies across entities and carries type-specific
information. The mean relation shift (difference in relation component
between R0 and R1) is 1.1, while the subject component has standard
deviation~0.24, giving a gating signal-to-noise ratio (SNR) of
approximately 4.7 for relation-specific neurons.

The pre-activation histograms (Figure~\ref{fig:preact}) directly
visualize this decomposition: two distributions of similar width
(reflecting the same subject variation), shifted apart by a constant
(the relation offset). GELU's zero crossing falls between them,
converting this soft shift into near-binary gating.

%----------------------------------------------------------------------
\subsection{Learned orthogonality eliminates interference}
\label{sec:orthogonality}
%----------------------------------------------------------------------

A natural concern is that the subject-varying signal might sometimes
flip a neuron's gate --- pushing a neuron that should be off into the
active regime, or vice versa. We find that this interference is
minimal, and the reason is geometric.

The relation direction in MLP input space (the difference between
R0 and R1 mean inputs) is \textbf{99\% orthogonal} to the subject
subspace (the top-10 principal components of subject variation; largest
cosine with any PC is 0.06). This near-orthogonality means that each
$\fc_1$ row has two nearly independent projections:
\begin{itemize}
\item \textbf{Relation-direction projection} ($\sim$25\% of row
  norm): determines gating. The sign perfectly predicts neuron
  ownership --- all 85 R0-specific neurons have positive projection,
  all 82 R1-specific neurons have negative projection, and dead
  neurons are near zero.
\item \textbf{Subject-subspace projection} ($\sim$48\% of row norm):
  determines type-varying activation magnitude within the active
  regime.
\end{itemize}
A single weight vector simultaneously implements both functions ---
gating and type computation --- without cross-talk, because the two
signals occupy orthogonal subspaces. The actual interference rate is
correspondingly low: only 3\% of (neuron, entity) pairs show
meaningful wrong-relation activation (GELU output $> 0.05$).

This orthogonality is not built into the architecture. It is a
\emph{learned} property: the model has discovered that placing the
relation signal in a direction orthogonal to subject variation is the
cleanest way to implement relation-conditioned gating without
corrupting the type-specific computation. Full per-neuron statistics
are in Appendix~\ref{app:neuron_decomp}.

%----------------------------------------------------------------------
\subsection{Connection to the gating mechanism}
\label{sec:attn_gating_connection}
%----------------------------------------------------------------------

The pre-activation shift of~1.1 between matching and non-matching
neurons (Section~\ref{sec:gelu_gating}) is produced by head~1's
contribution to the residual stream. Head~1 injects a
relation-identity vector $\bv_{\text{rel}}$; $\fc_1$ projects this
vector onto each neuron's weight row. The sign and magnitude of this
projection determine whether the neuron crosses GELU's zero point:
\begin{equation}
\fc_{1,j} \cdot \bv_{\text{R0}} - \fc_{1,j} \cdot \bv_{\text{R1}}
  = \fc_{1,j} \cdot (\bv_{\text{R0}} - \bv_{\text{R1}})
  \;\propto\; \text{gating shift for neuron } j.
\label{eq:gating_shift}
\end{equation}
The correlation between this $\fc_1$ projection and the empirically
measured gating shift is 1.000, confirming that the decomposition is
exact.

%----------------------------------------------------------------------
\subsection{Information availability before attention}
\label{sec:h0_info}
%----------------------------------------------------------------------

That separate affine maps achieve 99.7\% Voronoi accuracy from $\bh_0$
alone (Table~\ref{tab:residual}) demonstrates that the 64-dimensional
embedding already encodes both f-type and g-type structure in a
linearly accessible form. However, at inference the model has a
\emph{single} MLP with one set of weights --- it cannot condition on
the relation without reading the relation token. The attention output
is what enables this: by adding a relation-identity direction to the
residual stream, it allows one MLP to implement two different
transformations.

%======================================================================
\section{GELU Gating: Two Sub-Networks from One MLP}
\label{sec:gelu_gating}
%======================================================================

%----------------------------------------------------------------------
\subsection{Neuron partition}
\label{sec:partition}
%----------------------------------------------------------------------

Classifying each of the 256 MLP neurons by their mean post-GELU
activation (threshold 0.05) yields a clean partition:
\begin{center}
\begin{tabular}{lc}
\toprule
Category & Count \\
\midrule
R0-only & 85 \\
R1-only & 82 \\
Shared & 7 \\
Dead (both $< 0.05$) & 82 \\
\bottomrule
\end{tabular}
\end{center}
The MLP is effectively two independent sub-networks that happen to
share weights, routed by GELU gating. Figure~\ref{fig:activations}
shows the per-entity activation distributions for relation-specific
neurons: matching-relation activations spread broadly from 0 to~2,
while wrong-relation activations cluster tightly at GELU's minimum
($\approx -0.17$).

\begin{figure}[t]
\centering
\includegraphics[width=\textwidth]{results/neuron_activation_histograms.png}
\caption{Post-GELU activations at relation-specific neurons. Blue:
  matching relation context. Coral: wrong relation context. 2025
  entities $\times$ neurons per panel.}
\label{fig:activations}
\end{figure}

%----------------------------------------------------------------------
\subsection{Pre-activation analysis}
\label{sec:preact}
%----------------------------------------------------------------------

The gating mechanism is visible in the pre-GELU values
(Figure~\ref{fig:preact}):

\begin{table}[h]
\centering
\begin{tabular}{lcc}
\toprule
Context $\to$ Neurons & Mean & Std \\
\midrule
R0 contexts at R0 neurons & $+0.44$ & 0.34 \\
R1 contexts at R0 neurons & $-0.70$ & 0.43 \\
R1 contexts at R1 neurons & $+0.46$ & 0.39 \\
R0 contexts at R1 neurons & $-0.69$ & 0.45 \\
\bottomrule
\end{tabular}
\end{table}

\noindent The matching-relation distribution (mean $\approx +0.45$) and
wrong-relation distribution (mean $\approx -0.69$) are separated by
approximately 1.1 in pre-activation space. Both are roughly Gaussian
with similar spread, consistent with the additive decomposition of
Eq.~\eqref{eq:pre_decomp}: a constant relation shift plus
entity-varying subject noise. GELU's zero crossing falls between the
two distributions, converting this moderate shift into near-binary
on/off behavior.

\begin{figure}[t]
\centering
\includegraphics[width=\textwidth]{results/neuron_preactivation_histograms.png}
\caption{Pre-GELU values at relation-specific neurons. The two
  distributions have similar width (subject variation) but different
  means (relation offset). GELU's zero crossing falls between them.}
\label{fig:preact}
\end{figure}

%----------------------------------------------------------------------
\subsection{GELU derivative analysis}
\label{sec:gelu_derivative}
%----------------------------------------------------------------------

To quantify how linear GELU is in its operating regime, we evaluate
$\relu'(x)$ at each (entity, neuron) pre-activation value
(Figure~\ref{fig:derivative}):

\begin{table}[h]
\centering
\begin{tabular}{lcc}
\toprule
Context $\to$ Neurons & Mean $\relu'$ & Std \\
\midrule
Matching relation & 0.80 & 0.20 \\
Wrong relation & 0.09 & 0.21 \\
\bottomrule
\end{tabular}
\end{table}

\noindent Active neurons operate in GELU's near-linear regime
(derivative $\approx 0.8$, close to~1), while gated-off neurons are
in the saturated regime (derivative $\approx 0.09$, close to~0). The
distribution is bimodal with few values in the 0.2--0.5 transition
zone. GELU acts primarily as a \textbf{soft binary switch} rather than
providing graded nonlinearity.

The second derivative $\relu''(x) = \phi(x)(2 - x^2)$ is
non-negligible (mean 0.45--0.61), indicating local curvature. However,
this curvature does not accumulate into large deviations from
linearity: the affine approximation still achieves $R^2 \approx 0.98$
(Section~\ref{sec:affine}). The entity-specific variation in
pre-activations (std $\sim 0.3$--$0.4$) is small enough that the
nonlinear correction term remains minor.

\begin{figure}[t]
\centering
\includegraphics[width=\textwidth]{results/neuron_derivative_histograms.png}
\caption{GELU derivative $\relu'(x)$ at pre-activation values.
  Bimodal: near~0 (gated off) or near 0.8--1.0 (active). Few values
  in the transition zone.}
\label{fig:derivative}
\end{figure}

%----------------------------------------------------------------------
\subsection{Summary: the division of labor}
\label{sec:division}
%----------------------------------------------------------------------

The complete mechanism for implementing relation-conditioned affine
maps is:

\begin{enumerate}
\item \textbf{Attention (head~1)} adds a relation-identity direction
  $\bv_{\text{rel}}$ to the residual stream, orthogonal to the subject
  embedding subspace (Section~\ref{sec:superposition}).

\item \textbf{$\fc_1$} projects the superposed signal. The relation
  component determines gating; the subject component provides
  type-varying activation. These two functions coexist without
  interference due to the learned orthogonality
  (Section~\ref{sec:orthogonality}).

\item \textbf{GELU} converts the soft, relation-dependent offset into
  near-binary on/off gating, partitioning the 256 neurons into
  relation-specific sub-networks.

\item \textbf{$\fc_2$} maps the sparse, relation-specific activation
  vector to the 64-dimensional output space via a learned dictionary
  (Section~\ref{sec:dictionary}).
\end{enumerate}

\noindent The net effect is
$\mathrm{MLP}(\bx) \approx \bW_k \bx + \bb_k$,
where $k$ indexes the relation.

%======================================================================
\section{The $\fc_2$ Dictionary}
\label{sec:dictionary}
%======================================================================

The output weight matrix $\fc_2 \in \R^{64 \times 256}$ acts as a
dictionary: the MLP output is a weighted sum of its 256 column vectors,
where the weights are the post-GELU activations. The dictionary has
several notable properties:

\paragraph{Approximate isotropy.}
Mean pairwise cosine between columns is $\approx 0.001$ with condition
number 4.83 --- the columns spread roughly uniformly across $\R^{64}$,
providing a well-conditioned basis for representing arbitrary
directions through appropriate weighting.

\paragraph{Relation partition.}
The column partition mirrors the neuron partition: 85 R0-only, 82
R1-only, 7 shared, 82 dead. Each relation's sub-network uses a
distinct subset of dictionary elements.

\paragraph{Learned directions.}
The columns are not inherited from initialization. The mean cosine
between initial and trained directions is 0.67 --- every column
rotates substantially. Functionally important neurons rotate
\emph{more} (mean cosine 0.57) and grow more in norm ($+87\%$ vs.\
$+52\%$ average). The dictionary is shaped by training, not random
initialization.

\paragraph{Rank constraint.}
Because the MLP input is 64-dimensional, the 256-dimensional
activation vector is constrained to a 64-dimensional affine subspace
(the image of $\relu \circ \fc_1$). The network cannot produce
arbitrary activation patterns; it can only produce those reachable
from a 64-dimensional input through the gated nonlinearity.

%======================================================================
\section{Type-Level Computation}
\label{sec:type_level}
%======================================================================

The model's task structure implies that subjects sharing a type should
receive similar MLP computations. We quantify this by comparing
per-subject neuron contribution profiles (the 256-dimensional vector
of each neuron's signed contribution to the output).

\begin{table}[h]
\centering
\begin{tabular}{lcc}
\toprule
Metric & R0 & R1 \\
\midrule
Within-type contribution cosine & 0.911 & 0.917 \\
Across-type contribution cosine & 0.474 & 0.272 \\
Between-type / within-type variance & 221$\times$ & 363$\times$ \\
\bottomrule
\end{tabular}
\end{table}

\noindent Subjects of the same type receive nearly identical MLP
computations (cosine 0.91), while different types have substantially
different profiles. The between-type variance exceeds within-type
variance by 221--363$\times$. The MLP implements an approximately
\emph{type-uniform} computation.

The GELU gating structure reinforces this: 63--65\% of neurons are
effectively dead (output $< 0.01$) for any given input, and 113--116
neurons have perfect within-type consistency in their on/off state.
Different types activate different neuron subsets, implementing
type-conditioned computation through sparse gating.

The residual within-type variation ($\sim$9\% in cosine) is small but
nonzero. These 79--87 neurons with within-type gating variability
likely serve the harder object-selection problem: choosing a specific
object from among the 45 candidates sharing the target type, where
per-subject identity matters beyond type membership.

%======================================================================
\section{Discussion}
\label{sec:discussion}
%======================================================================

We have provided a complete mechanistic account of how a one-layer
transformer implements relation-conditioned affine transformations.
The mechanism has three steps, each cleanly attributable to a specific
architectural component:

\begin{enumerate}
\item A single attention head (head~1) reads the relation token and
  injects a relation-identity direction into the residual stream. This
  direction is superimposed on the subject embedding and 99\%
  orthogonal to it.

\item $\fc_1$ projects the combined signal. Because of the learned
  orthogonality, the relation projection (gating) and subject
  projection (type computation) are independent. The relation
  component shifts each neuron's pre-activation by $\sim$1.1,
  crossing GELU's threshold to produce a clean 85/82 neuron partition.

\item The active neurons' $\fc_2$ columns form a relation-specific
  dictionary that maps the subject-dependent activation variation to
  object type centroids.
\end{enumerate}

\paragraph{Linearity as an emergent property.}
The affine map is not a design constraint --- it emerges from
training. GELU is a nonlinear activation function, yet the trained
model uses it primarily as a binary gate. Within each relation's
active sub-network, the computation is approximately linear ($R^2
\approx 0.98$). The nonlinearity is concentrated at the gating
boundary (the transition between on and off), not within the active
computation. This suggests that when the task structure admits a
linear solution, gradient descent finds it, using the available
nonlinearity only where strictly necessary (for relation selection).

\paragraph{Learned orthogonality as the enabling geometry.}
The key geometric property is the 99\% orthogonality between the
relation direction and the subject subspace. Without this, the
gating and computation functions would interfere: subject variation
would flip neuron gates, and the relation signal would corrupt
type-specific activations. The model learns to place these signals
in orthogonal subspaces, allowing a single $\fc_1$ weight matrix to
serve dual purposes. This is a minimal instance of the more general
phenomenon of \emph{feature superposition} in transformer residual
streams \citep{elhage2022superposition}, where multiple independent
features coexist in a shared representational space by occupying
nearly orthogonal directions.

\paragraph{Limitations.}
Our analysis is specific to a one-layer model on a synthetic task
with two relations and 45 types. The mechanism's reliance on a single
attention head for relation routing and a clean neuron partition
reflects the simplicity of the setting. In deeper models or natural
language, the same functional decomposition (routing $\to$ gating
$\to$ readout) may be distributed across layers, with more complex
superposition structures and softer neuron partitions.

\paragraph{Future directions.}
Key questions include whether the superposition mechanism scales to
more relations (requiring more than a single direction), whether
multi-layer models distribute the mechanism across layers, and whether
analogous orthogonality structures can be identified in language
models trained on natural data.

%======================================================================
\appendix
\section{Detailed Attention Decomposition}
\label{app:attention}
%======================================================================

We decompose the attention output at position~2 into per-head
contributions. Since the output projection is linear, $\attnout =
\sum_{h=0}^{3} \attnout^{(h)}$, where each $\attnout^{(h)}$ is
head~$h$'s weighted values passed through the corresponding slice of
the output projection matrix (decomposition error $< 10^{-7}$).

\paragraph{Attention weight patterns.}
Table~\ref{tab:attn_weights} shows the mean attention weights at
position~2 attending to positions~0 (\texttt{sep}), 1 ($r_k$), and 2
(self). The patterns are stable across subjects (std 0.02--0.08).

\begin{table}[h]
\centering
\caption{Mean attention weights at position~2 $\to$ [sep, rel, self].}
\label{tab:attn_weights}
\begin{tabular}{ccccccc}
\toprule
& \multicolumn{3}{c}{R0 context} & \multicolumn{3}{c}{R1 context} \\
\cmidrule(lr){2-4} \cmidrule(lr){5-7}
Head & $\to$sep & $\to$rel & $\to$self & $\to$sep & $\to$rel & $\to$self \\
\midrule
0 & 0.743 & 0.146 & 0.110 & 0.748 & 0.142 & 0.111 \\
1 & 0.217 & 0.728 & 0.055 & 0.180 & 0.775 & 0.046 \\
2 & 0.193 & 0.700 & 0.106 & 0.167 & 0.739 & 0.093 \\
3 & 0.391 & 0.226 & 0.383 & 0.414 & 0.181 & 0.405 \\
\bottomrule
\end{tabular}
\end{table}

Head~1 and head~2 attend primarily to the relation token (70--78\%).
Head~0 attends primarily to \texttt{sep} (74\%). Head~3 splits between
\texttt{sep} and self.

\paragraph{Per-head output statistics.}
Table~\ref{tab:per_head} summarizes the output of each head.

\begin{table}[h]
\centering
\caption{Per-head attention output analysis. ``Between-rel var'' is the
  fraction of output variance attributable to relation identity.
  ``Within-rel cos'' is mean pairwise cosine of outputs under the same
  relation. ``Cross-rel cos'' is cosine between same-subject outputs
  under R0 vs.\ R1.}
\label{tab:per_head}
\begin{tabular}{cccccc}
\toprule
Head & R0 norm & R1 norm & $\cosim(\bar{\mu}_{\text{R0}}, \bar{\mu}_{\text{R1}})$
  & Between-rel var & Within-rel cos \\
\midrule
0 & 0.127 & 0.148 & $+0.988$ & 63\% & 0.995 \\
1 & 0.307 & 0.256 & $-0.481$ & 99.2\% & 0.999 \\
2 & 0.140 & 0.158 & $+0.656$ & 91\% & 0.996 \\
3 & 0.074 & 0.056 & $+0.242$ & 50\% & 0.78 \\
\bottomrule
\end{tabular}
\end{table}

Head~1 dominates: 99.2\% of its variance is between relations, its
outputs are near-identical within a relation (cosine 0.999), and its
R0/R1 means point in opposite directions ($\cosim = -0.48$). Its norm
(0.26--0.31) is 2--4$\times$ that of the other heads.

\paragraph{Relation direction decomposition.}
The total relation direction (mean$_{\text{R0}} -
\text{mean}_{\text{R1}}$, norm = 0.489) decomposes across heads:

\begin{center}
\begin{tabular}{ccc}
\toprule
Head & Projection onto rel.\ direction & \% of total \\
\midrule
0 & $-0.013$ & $-2.7\%$ \\
1 & $+0.465$ & $+95.2\%$ \\
2 & $+0.023$ & $+4.8\%$ \\
3 & $+0.014$ & $+2.8\%$ \\
\bottomrule
\end{tabular}
\end{center}

\paragraph{Subject-specific residuals.}
After subtracting each relation's mean output, the residual
per-head outputs carry subject-specific information:

\begin{center}
\begin{tabular}{cccc}
\toprule
Head & Subject var (R0) & Subject var (R1) & Cross-rel residual cos \\
\midrule
0 & 0.0001 & 0.0001 & 0.74 \\
1 & 0.0005 & 0.0005 & 0.26 \\
2 & 0.0003 & 0.0005 & 0.78 \\
3 & 0.0008 & 0.0025 & 0.79 \\
\bottomrule
\end{tabular}
\end{center}

\noindent Subject variance is small for all heads ($< 0.003$).
Heads~0, 2, 3 show moderate cross-relation consistency (0.74--0.79) in
their subject residuals, meaning the subject-specific component is
shared across relations. Head~1's subject residual is nearly
uncorrelated across relations (0.26) --- it is specialized for
relation encoding.

%======================================================================
\section{Per-Neuron Input Decomposition}
\label{app:neuron_decomp}
%======================================================================

We decompose each neuron's pre-activation into a relation component
(constant per relation) and a subject component (varying across
entities), as described in Eq.~\eqref{eq:pre_decomp}. The
decomposition is exact (max error $9.5 \times 10^{-7}$).

\paragraph{Gating SNR by neuron category.}

\begin{table}[h]
\centering
\caption{Per-neuron gating statistics by category. Gating shift =
  $|\text{rel\_comp}_{\text{R0}} - \text{rel\_comp}_{\text{R1}}|$.
  SNR = gating shift / subject std.}
\label{tab:gating_snr}
\begin{tabular}{lccccc}
\toprule
Category & $n$ & Mean shift & Subject std & Mean SNR & Median SNR \\
\midrule
R0-only & 85 & 1.14 & 0.25 & 4.71 & 4.90 \\
R1-only & 82 & 1.15 & 0.27 & 4.46 & 4.54 \\
Dead & 82 & 0.02 & 0.23 & 1.63 & 1.55 \\
\midrule
All & 256 & 0.86 & 0.24 & 3.53 & 3.87 \\
\bottomrule
\end{tabular}
\end{table}

Relation-specific neurons have gating SNR $\approx 4.7$: the relation
shift is nearly 5$\times$ the subject variation. Dead neurons have
near-zero shift --- they are suppressed for both relations.

\paragraph{Gate interference rates.}

\begin{table}[h]
\centering
\caption{How often subject variation flips the intended gate state.}
\label{tab:interference}
\begin{tabular}{lc}
\toprule
Condition & Rate \\
\midrule
R0 neurons activated by R1 context (pre $> 0$) & 5.6\% \\
R1 neurons activated by R0 context (pre $> 0$) & 6.1\% \\
R0 neurons suppressed in R0 context (pre $< 0$) & 7.8\% \\
R1 neurons suppressed in R1 context (pre $< 0$) & 10.5\% \\
\midrule
R0 neurons with GELU output $> 0.05$ in R1 context & 2.9\% \\
R1 neurons with GELU output $> 0.05$ in R0 context & 3.3\% \\
\bottomrule
\end{tabular}
\end{table}

\noindent Raw interference (pre-activation crossing zero) occurs in
6--10\% of cases, but GELU's nonlinearity near zero suppresses most of
these: only 3\% produce meaningful activation ($> 0.05$).

\paragraph{$\fc_1$ row geometry.}
Each $\fc_1$ row (dim=64) can be decomposed into its projection onto
the relation direction and the subject subspace
(Table~\ref{tab:fc1_geometry}). The correlation between the $\fc_1$
relation-direction projection and the empirically measured gating shift
is 1.000.

\begin{table}[h]
\centering
\caption{$\fc_1$ row projections by neuron category. ``Rel.\ dir.''
  is the signed projection onto the relation direction. All R0 neurons
  project positively; all R1 neurons negatively.}
\label{tab:fc1_geometry}
\begin{tabular}{lcccc}
\toprule
Category & $\fc_1 \cdot \hat{\mathbf{d}}_{\text{rel}}$ & $\|\fc_1\|_{\text{subj}}$ & $\|\fc_1\|$ & Rel.\ frac. \\
\midrule
R0-only & $+0.101 \pm 0.042$ & $0.143 \pm 0.058$ & 0.32 & 31\% \\
R1-only & $-0.103 \pm 0.050$ & $0.162 \pm 0.062$ & 0.34 & 30\% \\
Dead & $+0.002 \pm 0.039$ & $0.135 \pm 0.046$ & 0.26 & 1\% \\
\bottomrule
\end{tabular}
\end{table}

\paragraph{Example neurons.}
Table~\ref{tab:example_neurons} shows the complete decomposition for
representative neurons.

\begin{table}[h]
\centering
\caption{Example neuron profiles. ``Rel.\ comp.'' is the constant
  pre-activation component for each relation.}
\label{tab:example_neurons}
\begin{tabular}{lccccccc}
\toprule
 & \multicolumn{2}{c}{$\fc_1$ projections} & & \multicolumn{2}{c}{Rel.\ component} & \\
\cmidrule(lr){2-3} \cmidrule(lr){5-6}
Neuron & Rel.\ dir. & Subj.\ space & SNR & R0 & R1 & Role \\
\midrule
154 & $+0.186$ & 0.203 & 6.1 & $+0.78$ & $-1.30$ & R0-specific \\
157 & $-0.296$ & 0.290 & 7.5 & $-1.77$ & $+1.56$ & R1-specific \\
15  & $+0.001$ & 0.167 & 0.04 & $-0.93$ & $-0.94$ & Dead \\
\bottomrule
\end{tabular}
\end{table}

\noindent Neuron~154 (top R0): nearly half its $\fc_1$ norm is in the
relation direction, giving a gating shift of 2.08 (SNR 6.1). Under R0,
it operates well into GELU's linear regime (mean pre-activation
$+0.78$); under R1, deeply suppressed ($-1.30$). Neuron~157 (top R1)
is the mirror image with even higher SNR (7.5). Neuron~15 (lowest SNR)
has essentially zero relation-direction projection and is suppressed
for both relations (pre-activation $\approx -0.93$).

%======================================================================
\bibliography{references}
%======================================================================

\end{document}
